\bookchapter{Chapitre 7 : Analyse complète des fonctions de pointe}

\chapterintro{La technologie évolue sans cesse, et les fonctionnalités du smartphone aussi. Communication par satellite, communication hors réseau, détection d'accident, technologie SparkLink... Ces fonctions qui semblent sorties de la science-fiction ont en fait déjà intégré notre quotidien. Découvrons ensemble ces technologies de pointe.}

\booksection{Communication par satellite : demander de l'aide même sans signal}

La communication par satellite est une avancée majeure dans le domaine des smartphones ces dernières années. Quand vous êtes en pleine nature sauvage, dans les montagnes profondes, sans aucun signal mobile, la communication par satellite peut être votre bouée de sauvetage.

\booksubsection{Principe de fonctionnement :}

Le smartphone établit directement une connexion avec un satellite de communication en orbite pour envoyer un message de détresse. Actuellement, il prend principalement en charge l'envoi de SMS et d'informations de localisation, pas les appels téléphoniques.

\booksubsection{Modèles pris en charge :}

\begin{itemize}
  \item \textbf{Apple}: À partir de la série iPhone 14, l'envoi d'un SOS d'urgence par satellite est pris en charge. La série iPhone 17 permet en plus d'envoyer des iMessage et des SMS par satellite, avec une couverture élargie à davantage de régions.
  \item \textbf{Huawei}: La série Mate 50 commence à prendre en charge les messages par satellite Beidou (texte uniquement), la série Mate 60 Pro prend en plus en charge les appels par satellite Tiantong (possibilité d'appeler), la série Mate 70 renforce encore les capacités de communication par satellite.
  \item \textbf{Xiaomi/OPPO/vivo}: Depuis 2024, les modèles phares prennent progressivement en charge la fonction de messages par satellite Beidou.
  \item \textbf{Tendance de développement}: La communication par satellite passe d'une fonction exclusive aux modèles phares haut de gamme à une diffusion vers les modèles intermédiaires. Le système de satellites Beidou chinois bénéficie de la prise en charge la plus large sur les smartphones nationaux.
\end{itemize}

\booksubsection{Scénarios d'utilisation :}

\begin{itemize}
  \item Randonnée aventureuse, détresse en montagne.
  \item Détresse en mer.
  \item Catastrophes naturelles coupant les communications terrestres.
\end{itemize}

\booksubsection{Méthode d'utilisation :}

\begin{enumerate}
  \item Assurez-vous d'être dans un espace dégagé, avec une vue sur le ciel.
  \item Ouvrez la fonction de communication par satellite du smartphone (généralement dans l'interface d'appel d'urgence).
  \item Suivez les instructions à l'écran pour pointer vers le satellite.
  \item Envoyez le message de détresse prédéfini et votre localisation.
\end{enumerate}

\begin{tcbinfo}{Le saviez-vous}
La communication par satellite nécessite une antenne matérielle spéciale dans le smartphone, tous les smartphones ne la prennent pas en charge. De plus, la communication par satellite nécessite généralement un abonnement payant, mais la fonction de détresse d'urgence est généralement gratuite.

\end{tcbinfo}

\booksection{Communication hors réseau : un nouveau mode de communication à courte distance}

La « communication hors réseau » désigne la technologie permettant à deux smartphones de communiquer directement entre eux en l'absence de réseau mobile et de Wi-Fi.

\booksubsection{Principe technique :}

Utilise le Bluetooth ou des fréquences radio dédiées pour établir une connexion point à point à courte distance (généralement de quelques centaines de mètres à quelques kilomètres).

\booksubsection{Scénarios d'application :}

\begin{itemize}
  \item \textbf{Activités de plein air}: Communication entre membres d'une équipe en randonnée ou en camping.
  \item \textbf{Secours en cas de catastrophe}: Quand les communications sont coupées après un séisme ou une inondation.
  \item \textbf{Grands événements}: Concerts, événements sportifs avec une forte affluence et un signal faible.
\end{itemize}

\booksubsection{Plateformes prises en charge :}

\begin{itemize}
  \item \textbf{Apple}: « NameDrop » et « Trouver des amis hors ligne » d'iOS 17.
  \item \textbf{Huawei}: L'application Canghai prend en charge les appels hors réseau.
  \item \textbf{Xiaomi}: Certaines fonctions du centre Xiaomi Miaoxiang.
\end{itemize}

\booksection{Détection d'accident de voiture : sauver des vies au moment critique}

La détection d'accident est une fonctionnalité qui utilise les capteurs du smartphone pour détecter un accident grave et appeler automatiquement les secours.

\booksubsection{Principe de fonctionnement :}

Le smartphone utilise plusieurs capteurs — accéléromètre, gyroscope, baromètre, micro — pour déterminer de manière globale si une collision grave s'est produite. Après détection, le smartphone :

\begin{enumerate}
  \item Affiche un compte à rebours, vous demandant si vous allez bien.
  \item Si vous ne répondez pas à la fin du compte à rebours, appelle automatiquement les secours.
  \item Envoie votre localisation à vos contacts d'urgence.
\end{enumerate}

\booksubsection{Modèles pris en charge :}

\begin{itemize}
  \item \textbf{Apple}: iPhone 14 et ultérieurs.
  \item \textbf{Android}: Certains modèles haut de gamme prennent progressivement en charge cette fonction.
\end{itemize}

\booksubsection{Précautions :}

\begin{itemize}
  \item La fonction doit être activée pour être utilisée.
  \item Peut se déclencher par erreur (comme sur des montagnes russes ou lors d'un exercice intense), il suffit d'annuler à temps.
  \item Dépend des capteurs et des algorithmes, pas fiable à 100 \%, mais constitue une couche de protection supplémentaire.
\end{itemize}

\booksection{Technologie SparkLink : la connexion sans fil de nouvelle génération}

SparkLink est une technologie de communication sans fil de nouvelle génération développée sous l'impulsion de Huawei, visant à remplacer le Bluetooth et le Wi-Fi dans certains scénarios. Elle surpasse les solutions existantes sur les indicateurs clés : débit en bande de base jusqu'à 900 Mbps (largement supérieur aux 50 Mbps du Bluetooth), latence aussi basse que 20 $\mu$ s (Bluetooth environ 15 ms), réseau unique prenant en charge 4096 appareils connectés simultanément (Bluetooth généralement moins de 8).

\booksubsection{Caractéristiques techniques :}

\begin{itemize}
  \item \textbf{Ultra-faible latence}: Latence de bout en bout de seulement 20 microsecondes (0,02 ms), deux ordres de grandeur inférieurs au Bluetooth, la synchronisation jeu et musique est quasi imperceptible.
  \item \textbf{Débit élevé}: Débit maximal par liaison de 900 Mbps, des dizaines de fois supérieur à l'audio Bluetooth, permet de transmettre de l'audio sans perte voire de la vidéo HD.
  \item \textbf{Faible consommation}: Consommation en veille comparable au Bluetooth, mais efficacité de transmission supérieure, énergie par donnée plus faible.
  \item \textbf{Couverture étendue + haute concurrence}: Portée de connexion supérieure de plus de 50 \% au Bluetooth, réseau unique prenant en charge plus de mille appareils simultanément.
\end{itemize}

\booksubsection{Scénarios d'application :}

\begin{itemize}
  \item \textbf{Écouteurs sans fil}: Latence plus faible, synchronisation audio-vidéo pour les jeux.
  \item \textbf{Montres connectées}: Connexion plus stable.
  \item \textbf{Maison connectée}: Plus d'appareils connectés simultanément.
  \item \textbf{Connectivité embarquée}: Connexion smartphone-système embarqué plus fluide.
\end{itemize}

\booksubsection{Appareils pris en charge :}

Actuellement principalement les appareils de l'écosystème Huawei, comme la série Huawei Mate 60, certains écouteurs et montres Huawei. D'autres marques prendront en charge à l'avenir.

\textbf{Perspectives technologiques}\par

La technologie SparkLink représente l'orientation future des communications sans fil. Avec l'adhésion de davantage de fabricants, nous verrons probablement davantage d'appareils prenant en charge SparkLink, pour une meilleure expérience sans fil.

\begin{tcbpractice}{Vérifiez : votre smartphone prend-il en charge ces fonctionnalités}
\begin{enumerate}
  \item Recherchez « satellite » dans les paramètres pour voir s'il y a des options de communication par satellite.
  \item Recherchez « détection d'accident » ou « SOS d'urgence », vérifiez si c'est activé.
  \item Si votre smartphone est un Huawei, vérifiez s'il prend en charge la connexion d'appareils SparkLink.
  \item Apprenez à utiliser ces fonctionnalités, elles pourraient vous servir au moment critique.
\end{enumerate}

\end{tcbpractice}
